\documentclass[aspectratio=169, 11pt]{beamer}

% --- Packages & Setup ---
\usepackage[utf8]{inputenc}
\usepackage[T1]{fontenc}
\usepackage{xcolor}
\usepackage{booktabs} 
\usepackage{graphicx}
\usepackage{tikz} 
\usepackage{array}
\usepackage{tabularx}
\usepackage{amsmath, amssymb}
\usetikzlibrary{calc}

% Branding Colors
\definecolor{amritaMaroon}{HTML}{A4123F} 

% Theme Setup
\usetheme{Madrid}
\usecolortheme[named=amritaMaroon]{structure}
\setbeamertemplate{navigation symbols}{}
\setbeamertemplate{footline}[frame number]

% Custom Column Type for wrapping text in tables
\newcolumntype{L}{>{\raggedright\arraybackslash}X}

\begin{document}

% --- SLIDE 1: TITLE SLIDE ---
\begin{frame}
    \begin{tikzpicture}[remember picture,overlay]
        \node[fill=amritaMaroon, text=white, font=\bfseries\footnotesize, anchor=north east, rounded corners=2pt, inner sep=5pt] 
        at ($(current page.north east) + (-0.2cm,-0.2cm)$) 
        {Type 1: Research Project};
    \end{tikzpicture}

    \centering
    \includegraphics[width=3.5cm]{logo.png} 
    \vspace{0.3cm}
    
    {\Large \textbf{\textcolor{amritaMaroon}{To develop an edge-based deep learning speech therapy system for autistic children}}} \\
    \vspace{0.2cm}
    {\normalsize \textbf{Team Number: 96}} \\ 
    \vspace{0.5cm}

    \begin{table}
        \centering
        \small 
        \renewcommand{\arraystretch}{1.2}
        \begin{tabular}{|c|c|l|}
            \hline
            \textbf{S.No} & \textbf{Register Number} & \textbf{Student Name} \\
            \hline
            1 & CB.SC.U4CSE23018 & Rohith Kumar D \\
            \hline
            2 & CB.SC.U4CSE23055 & T Venkataramana \\
            \hline
            3 & CB.SC.U4CSE23058 & VAB Jashwanth Reddy \\
            \hline
            4 & CB.SC.U4CSE23060 & C Kalyan Kumar Reddy \\
            \hline
            5 & CB.SC.U4CSE23446 & Hemanth Sholingaram \\
            \hline
        \end{tabular}
    \end{table}
    \vspace{0.3cm}
    
    Guide: \textbf{Dr. Venkataraman D} \\
    {\footnotesize Associate Professor, Department of CSE}
    
    \vspace{0.5cm}
    {\footnotesize \today}
\end{frame}

% --- SLIDE: PROBLEM, OBJECTIVES & NEED ---
\begin{frame}{Problem Statement, Objectives \& Need}
    \vspace{-2mm}
    \setbeamercolor{block title}{bg=amritaMaroon, fg=white}
    \begin{block}{\small Problem Statement}
        \small \centering
       Existing ASD speech therapy lacks objective biomarkers for reliable progress tracking[5,16]

Severe shortages of trained speech therapists and child psychiatrists limit access to timely intervention [23]

Cloud-based ASR systems fail to handle atypical regional pediatric speech and raise privacy concerns [1,22]
    \end{block}

    \vspace{2mm}

    \begin{columns}[T, totalwidth=\textwidth]
        \begin{column}{0.48\textwidth}
            \setbeamercolor{block title}{bg=black, fg=white}
            \begin{block}{\small Project Objectives}
                \scriptsize
                \begin{itemize}
                    \setlength{\itemsep}{5pt}
                    \item \textbf{Privacy-First Edge AI:} Engineer local biomarker extraction to ensure data sovereignty [25].
                    \item \textbf{Regional Optimization:} Adapt DL models for Tamil and Telugu phonetic nuances in ASD [15].
                    \item \textbf{Adaptive Learning:} Use RL to personalize therapy difficulty based on child progress [23].
                \end{itemize}
            \end{block}
        \end{column}

        \begin{column}{0.48\textwidth}
            \setbeamercolor{block title}{bg=amritaMaroon, fg=white}
            \begin{block}{\small Need of the System}
                \scriptsize
                \begin{itemize}
                    \setlength{\itemsep}{5pt}
                    \item \textbf{Address Scarcity:} Bridges the critical gap in therapist availability in rural South India [15].
                    \item \textbf{Zero-Trust Security:} Keeps sensitive pediatric voice data offline on edge hardware [25].
                    \item \textbf{Clinical Accuracy:} Replaces subjective observation with repeatable, data-driven metrics [4,24].
                \end{itemize}
            \end{block}
        \end{column}
    \end{columns}
\end{frame}

% --- SURVEY SET 1 ---
\begin{frame}{Literature Survey Set 1: Institutional Foundations \& Deep Acoustic Metrics}
    \tiny
    \centering
    \renewcommand{\arraystretch}{1.1}
    \setlength{\tabcolsep}{2pt}
    \begin{tabularx}{\textwidth}{|c|L|L|L|L|L|}
        \hline
        \textbf{Sl} & \textbf{Title} & \textbf{Detailed Description} & \textbf{Advantage} & \textbf{Disadvantage} & \textbf{Remarks} \\ \hline
        1 & The Taste of IPA (2024) [1] & Proposes CLAP-IPA, a multilingual phoneme-speech embedding model for cross-lingual keyword spotting. & Achieves strong zero-shot performance across diverse languages without retraining. & Requires extensive phonetic expertise for creating manual high-quality transcriptions. & Universal IPA symbols act as cross-lingual modeling units. \\ \hline
        2 & VSA Assessment (2013) [5] & Uses GMM and K-Means clustering to automate Vowel Space Area (VSA) estimates from connected speech. & More sensitive to speech production deficits than traditional corner-vowel methods. & Manual phonetic segmentation is cumbersome and often clinically unreliable. & VSA acts as a proxy for the kinematic displacement of articulators. \\ \hline
        3 & Digital Games (2022) [3] & Review of digital games focusing on engagement through interactive feedback and immersive environments. & Positive impact on satisfaction, motivation, and sustained attention in children. & Potential for frustration and low self-esteem after repeated in-game failures. & 48\% of designed games allow collaborative practice with parents. \\ \hline
        4 & AI Speech Class. (2025) [4] & Systematic review of AI architectures for 16 distinct speech and language disorders (SLDs). & Highlights the transition from traditional SVMs to hybrid deep learning (DL) architectures. & Lack of longitudinal and cross-cultural representative datasets for training. & Multimodal frameworks improve diagnostic accuracy significantly. \\ \hline
        5 & Emotion SLR (2022) [2] & Explores observation channels (facial, prosody, physiological) for ASD emotion detection. & Identifies early fusion as the most frequent multimodal integration method. & Channels are prone to environmental disturbance, causing data loss. & SVM and neural networks remain the most popular classifiers. \\ \hline
    \end{tabularx}
\end{frame}

% --- SURVEY SET 2 ---
\begin{frame}{Literature Survey Set 2: Scaffolding \& Scoped Interventions}
    \tiny
    \centering
    \renewcommand{\arraystretch}{1.1}
    \setlength{\tabcolsep}{2pt}
    \begin{tabularx}{\textwidth}{|c|L|L|L|L|L|}
        \hline
        \textbf{Sl} & \textbf{Title} & \textbf{Detailed Description} & \textbf{Advantage} & \textbf{Disadvantage} & \textbf{Remarks} \\ \hline
        6 & Tell-Your-Day AAC [6] & Visual AAC interface enabling autistic children to construct sentences about daily activities. & Increases sentence length and improves functional daily communication. & Limited templates require adult support for vocabulary expansion tasks. & Focuses on daily routine expressive communication scaffolding. \\ \hline
        7 & Robot Training [7] & Uses structured robot-guided interactions to improve verbal initiation and joint attention. & High engagement levels encourage spontaneous verbal responses in children. & Extremely high hardware cost and limited transfer to human interaction. & Technology-supported method for increasing verbal initiation. \\ \hline
        8 & AI Speech Feedback [8] & Real-time articulation and intelligibility feedback using ASR scoring for autistic speech. & Objective instant feedback enables frequent and independent speech practice. & ASR accuracy decreases significantly for highly atypical speech patterns. & Useful AI tool supporting clinician-guided articulation therapy. \\ \hline
        9 & Therapy Approaches [9] & Evaluation of traditional versus technology-enhanced speech therapies for ASD outcomes. & Combined approaches significantly improve speech and communication outcomes. & The individual contribution of each technology remains difficult to isolate. & Provides evidence for integrated multi-method therapy models. \\ \hline
        10 & ASR Assessment [10] & ASR-based systems designed to assess and enhance communication skills in children. & Enables objective speech evaluation and scalable digital therapy support. & Performance degrades with severely impaired or highly variable speech. & ASR effectively supports autism therapy as an assistive tool. \\ \hline
    \end{tabularx}
\end{frame}

% --- SURVEY SET 3 ---
\begin{frame}{Literature Survey Set 3: Scaffolding \& Virtual Interactions}
    \tiny
    \centering
    \renewcommand{\arraystretch}{1.1}
    \setlength{\tabcolsep}{2pt}
    \begin{tabularx}{\textwidth}{|c|L|L|L|L|L|}
        \hline
        \textbf{Sl} & \textbf{Title} & \textbf{Detailed Description} & \textbf{Advantage} & \textbf{Disadvantage} & \textbf{Remarks} \\ \hline
        11 & Immersive VR (2020) [11] & Creates safe virtual environments for practicing social scripts and public interactions. & Significantly lowers cognitive load for complex social tasks. & Can cause motion sickness, dizziness, and sensory overload in children. & Controlled environment for high-anxiety therapeutic tasks. \\ \hline
        12 & Visual Feedback (2021) [12] & Converts vocal pitch and volume into real-time graphics for visual profiling. & Provides immediate auditory-visual biofeedback for self-correction. & Risk of over-reliance on visual cues over phonetic accuracy. & Helps children visualize their own vocal performance. \\ \hline
        13 & Self-Learning Bot (2022) [13] & Implements "fading prompts" using Transformers for intent decoding in therapy. & Decodes intent effectively despite significant phonetic errors. & Requires massive datasets for fine-tuning Transformer architectures. & Simulates human therapist scaffolding and prompt fading. \\ \hline
        14 & Mobile Apps (2019) [14] & Replaces traditional cards with video modeling and adaptive articulation profiling. & Tracks phoneme clarity progress longitudinally for home practice. & Minimal validation in formal clinical settings for many apps. & Allows for accessible home-based articulation practice. \\ \hline
        15 & Knowledge Transfer (2023) [15] & Meta-learning approach to adapt global models to local dialects and accents. & Bridges the technology gap for Global South regional dialects. & High false-negative rates in cases of complex phonetic collisions. & Allows rapid deployment in regional and rural areas. \\ \hline
    \end{tabularx}
\end{frame}

% --- SURVEY SET 4 ---
\begin{frame}{Literature Survey Set 4: Prosody, Fusion \& Motor Learning}
    \tiny
    \centering
    \renewcommand{\arraystretch}{1.1}
    \setlength{\tabcolsep}{2pt}
    \begin{tabularx}{\textwidth}{|c|L|L|L|L|L|}
        \hline
        \textbf{Sl} & \textbf{Title} & \textbf{Detailed Description} & \textbf{Advantage} & \textbf{Disadvantage} & \textbf{Remarks} \\ \hline
        16 & Atypical Prosody (2021) [16] & openSMILE toolkit used to quantify robotic speech patterns via Jitter and Shimmer. & Established gold standard for extracting objective prosody metrics. & Hard to distinguish ASD-specific prosody from general speech delay. & Correlates with clinician perception of vocal monotony. \\ \hline
        17 & Multi-modal Fusion (2022) [17] & Syncs vocalization data with joint attention and gesture tracking systems. & Provides a comprehensive, holistic picture of child communication. & High setup complexity makes it difficult for home-based use. & Focuses on the synergy between gesture and speech. \\ \hline
        18 & Feedback Timing (2020) [18] & Study on the impact of immediate versus delayed feedback on motor learning. & Optimizes response timing for faster skill acquisition. & Rigid schedules do not adapt to the child's real-time mood. & Crucial for designing effective AI coaching interfaces. \\ \hline
        19 & Acoustic Modeling (2018) [19] & VTLN and GANs used to transform disordered audio into clear targets. & Compensates for pediatric vocal tract size and anatomy differences. & GAN-generated audio may sound unnatural or distorted. & Synthetic speech generation overcomes disordered data scarcity. \\ \hline
        20 & EPG Biofeedback (2017) [20] & Visualizes tongue-palate contact using specialized intra-oral sensors. & Highly effective at correcting stubborn misarticulation errors. & High hardware cost and significant physical discomfort. & Superior for visualizing specific misarticulation patterns. \\ \hline
    \end{tabularx}
\end{frame}

% --- SURVEY SET 5 ---
\begin{frame}{Literature Survey Set 5: Generative AI, Screening \& Edge Computing}
    \tiny
    \centering
    \renewcommand{\arraystretch}{1.4}
    \setlength{\tabcolsep}{3pt}
    \begin{tabularx}{\textwidth}{|c|L|L|L|L|L|}
        \hline
        \textbf{Sl} & \textbf{Title} & \textbf{Detailed Description} & \textbf{Advantage} & \textbf{Disadvantage} & \textbf{Remarks} \\ \hline
        21 & Neural Vocoders (2019) [21] & Compares WaveNet and SampleRNN for speech reconstruction from mel-spectrograms. & Outperforms traditional signal processing in speech naturalness. & Autoregressive models suffer from extremely slow generation speeds. & Establishes the trade-off between naturalness and speed. \\ \hline
        22 & Noor Project (2025) [22] & Fine-tunes Wav2Vec 2.0 Transformer models to classify ASD paralinguistics. & Achieves high accuracy (67\%) with scarce clinical data via transfer learning. & Struggles with class imbalance and performance drops across languages. & Proves Transformers capture paralinguistic biomarkers effectively. \\ \hline
        23 & DDPG Therapy (2025) [23] & Uses Deep Deterministic Policy Gradient (RL) to optimize screening strategies. & Achieves superior stability and 96.98\% accuracy over static models. & Deep RL requires massive interaction data to converge reliably. & Validates continuous control RL for dynamic intervention. \\ \hline
        24 & GenAI Therapy (2025) [24] & 6-week longitudinal study of Generative AI home therapy with 2000+ recordings. & Dose-response threshold (>15 hrs) led to significant gains in vocab. & Requires human-in-the-loop setup; not fully autonomous. & Provides clinical evidence for AI-boosted therapy outcomes. \\ \hline
        25 & Federated Privacy (2024) [25] & Introduces FLAMENCO: federated anomaly detection using Fully Homomorphic Encryption (FHE). & Uses FHE for rigorous privacy against malicious data attacks. & FHE encryption adds significant computation overhead and latency. & Proves gamified screening can be secure without centralization. \\ \hline
    \end{tabularx}
\end{frame}

% --- SLIDE: ORIGINALITY ---
\section{Originality}
\begin{frame}{ Originality of Work: A Hybrid Multimodal Approach}
    \begin{columns}
        \begin{column}{0.5\textwidth}
            \setbeamercolor{block title}{bg=amritaMaroon, fg=white}
            \begin{block}{The Novel Research Gap}
                Existing systems focus on cloud-based transcription [10]. Our study shifts to \textbf{Longitudinal Anomaly Analysis} using a privacy-first edge framework [25].
            \end{block}
            \textbf{Key Innovations:}
            \begin{itemize}
                \item \textbf{Privacy-Centric Federated TinyML:} Local on-device training to keep clinical voice data at home [25].
                \item \textbf{Cross-Modal Biomarkers:} Synchronizing spectral flux with non-intrusive visual stress markers [17].
            \end{itemize}
        \end{column}
        \begin{column}{0.5\textwidth}
            \begin{itemize}
                \item \textbf{RL Scaffolding:} Extending Seq2Seq models with "Fading Prompts" based on DQN motor stability [23].
                \item \textbf{Dravidian Normalization:} First adaptation of VTLN for South Indian retroflex consonants in ASD [19].
                \item \textbf{Anomaly-Driven Scoring:} Using OCSVM to distinguish therapeutic progress from emotional noise [4].
            \end{itemize}
        \end{column}
    \end{columns}
\end{frame}


% --- SLIDE 8: CUMULATIVE RESEARCH GAP ---
\section{Cumulative Gap}
\begin{frame}{Cumulative Research Gap [1--25]}
    \vspace{-4mm}
    \begin{columns}[T, totalwidth=\textwidth]
        
        % --- COLUMN 1: TECHNICAL & PRIVACY ---
        \begin{column}{0.48\textwidth}
            \setbeamercolor{block title}{bg=amritaMaroon, fg=white}
            \begin{block}{\small Technical \& Privacy Gaps}
                \scriptsize
                \begin{itemize}
                    \item \textbf{Privacy-Performance Trade-off:} Existing high-accuracy models rely on cloud-based APIs [10], while current secure Edge-AI/TinyML frameworks suffer from accuracy loss in disordered speech [25].
                    \item \textbf{Latency in Encryption:} Secure Federated Learning and Homomorphic Encryption (FHE) protocols currently introduce prohibitive computational overhead on low-power therapeutic hardware [25].
                \end{itemize}
            \end{block}
            
            \vspace{2mm}
            
            \setbeamercolor{block title}{bg=black, fg=white}
            \begin{block}{\small Linguistic Inequity}
                \scriptsize
                \begin{itemize}
                    \item \textbf{Phonetic Bias:} A critical lack of ASR optimization exists for the specific retroflex consonants and vowel centralization patterns found in \textbf{local autism kids} [15,22].
                \end{itemize}
            \end{block}
        \end{column}

        % --- COLUMN 2: CLINICAL & ADAPTIVITY ---
        \begin{column}{0.48\textwidth}
            \setbeamercolor{block title}{bg=black, fg=white}
            \begin{block}{\small Clinical \& Adaptive Gaps}
                \scriptsize
                \begin{itemize}
                    \item \textbf{Static Scaffolding:} Most digital tools follow rigid prompt schedules [6]; there is an absence of real-time \textbf{Reinforcement Learning (RL)} agents that adapt difficulty based on child stress [23].
                    \item \textbf{Single-Modality Tracking:} Current systems focus almost exclusively on audio [16], failing to fuse acoustic biomarkers (VSA) with non-intrusive visual indicators of sensory overload [2,17].
                \end{itemize}
            \end{block}
            
            \vspace{2mm}
            
            \setbeamercolor{block title}{bg=amritaMaroon, fg=white}
            \begin{block}{\small Objective Progress Tracking}
                \scriptsize
                \begin{itemize}
                    \item \textbf{Longitudinal Analysis:} Clinical tools lack automated longitudinal tracking that correlates phonetic "merging" (VSA) with long-term motor stability [4,24].
                \end{itemize}
            \end{block}
        \end{column}
    \end{columns}

    \vspace{4mm}
    \begin{beamercolorbox}[sep=4pt, center, rounded=true]{palette primary}
        \footnotesize \textbf{The Foundational Gap:} The lack of an integrated, local, and regional-aware system that adapts to the unique needs of local autism kids without compromising data sovereignty [1,5,9,14,25].
    \end{beamercolorbox}
\end{frame}

% --- SLIDE: TECHNICAL KNOWLEDGE ---
\section{Technical Knowledge}
\begin{frame}{Technical Knowledge: Algorithmic Fundamentals}
    \small
    \begin{columns}[t]
        % --- Column 1 ---
        \begin{column}{0.48\textwidth}
            \setbeamercolor{block title}{bg=amritaMaroon, fg=white}
            
            \begin{block}{1. Acoustic Feature Pipeline}
                \begin{itemize}
                    \item \textbf{Spectral (MFCC):} Maps vocal tract shape to human auditory perception.
                    \item \textbf{Non-Linear (EMD):} Decomposes noisy signals to isolate vocal tremors.
                \end{itemize}
            \end{block}

            \vspace{0.5em}

            \begin{block}{2. Temporal Alignment}
                \begin{itemize}
                    \item \textbf{DTW Logic:} Minimizes time-distance cost to compare disordered speech against healthy templates.
                    \item Corrects for variations in speaking rate.
                \end{itemize}
            \end{block}
        \end{column}

        % --- Column 2 ---
        \begin{column}{0.48\textwidth}
            \setbeamercolor{block title}{bg=black, fg=white}
            
            \begin{block}{3. Anomaly Detection}
                \begin{itemize}
                    \item \textbf{OC-SVM:} Learns "normative" speech boundaries to flag deviations.
                    \item \textbf{GOP Score:} Measures pronunciation quality based on hyperplane distance.
                \end{itemize}
            \end{block}

            \vspace{0.5em}

            \setbeamercolor{block title}{bg=amritaMaroon, fg=white}
            \begin{block}{4. Articulatory Precision}
                \begin{itemize}
                    \item \textbf{Vowel Space (VSA):} Uses F1 and F2 formants to track tongue/jaw range.
                    \item \textbf{Clarity Metric:} Area reduction indicates slurred speech or "vowel centralization."
                \end{itemize}
            \end{block}
        \end{column}
    \end{columns}
\end{frame}
\begin{frame}{System Logic: Cross-Modal \& Adaptive Flow}
    \begin{columns}[T]
        \begin{column}{0.48\textwidth}
            \setbeamercolor{block title}{bg=amritaMaroon, fg=white}
            \begin{block}{Acoustic Biomarker Extraction}
                \scriptsize
                Uses \textbf{Vowel Space Area (VSA)} to track speech motor stability from formant patterns.
                \textbf{Spectral Analysis:} MFCC features help normalize and compare disordered speech.
            \end{block}
            
            \begin{block}{Temporal Alignment (DTW)}
                \scriptsize
                DTW aligns child speech with reference templates even when speaking speed changes.
                This improves fair pronunciation comparison over time.
            \end{block}
        \end{column}

        \begin{column}{0.48\textwidth}
            \setbeamercolor{block title}{bg=black, fg=white}
            \begin{block}{RL-Based Adaptive Scaffolding}
                \scriptsize
                Uses \textbf{DDPG (Deep Deterministic Policy Gradient)} to adjust therapy difficulty in real time:
                \begin{itemize}
                    \item \textbf{State:} Phonetic accuracy + Visual stress markers.
                    \item \textbf{Action:} Dynamic "Fading Prompt" intensity.
                    \item \textbf{Goal:} Balance speech accuracy and willingness.
                \end{itemize}
            \end{block}

            \begin{block}{Privacy-First Federated Learning}
                \scriptsize
                Uses \textbf{Fully Homomorphic Encryption (FHE)} for secure local anomaly detection.
                Keeps pediatric biometric voice data private on local systems.
            \end{block}
        \end{column}
    \end{columns}
\end{frame}

% --- SLIDE 3: TECHNICAL & RESEARCH CHALLENGES ---
\section{Challenges}
\begin{frame}{Technical \& Research Challenges}
\vspace{-4mm}
\begin{columns}[T, totalwidth=\textwidth]
\begin{column}{0.48\textwidth}
\setbeamercolor{block title}{bg=amritaMaroon, fg=white}
\begin{block}{\small Technical Implementation Challenges}
\scriptsize
\begin{itemize}
    \item \textbf{Edge Optimization:} Quantizing heavy Transformer architectures for ARM-based edge hardware [22,25].
    \item \textbf{Real-Time Multimodal Sync:} Zero-latency synchronization between audio and visual sensors [17].
    \item \textbf{Privacy:} Implementing Federated Learning without transmitting raw pediatric voice data [25].
\end{itemize}
\end{block}
\end{column}

\hspace{0.02\textwidth}
\begin{column}{0.48\textwidth}
\setbeamercolor{block title}{bg=black, fg=white}
\begin{block}{\small Research \& Scientific Challenges}
\scriptsize
\begin{itemize}
    \item \textbf{Data Scarcity:} Overcoming the lack of high-quality Dravidian disordered pediatric speech datasets [15,22].
    \item \textbf{Acoustic Modeling:} Distinguishing between vowel centralization and general speech delays [5,16].
    \item \textbf{RL Reward Engineering:} Defining reward functions that balance accuracy with emotional willingness [23].
\end{itemize}
\end{block}
\end{column}
\end{columns}
\end{frame}

% --- NEW SLIDE: ARCHITECTURAL PILLARS & MATHEMATICAL FOUNDATION ---


% --- NEW SLIDE: REFINED PROBLEM STATEMENT ---
\begin{frame}{Refined Problem Statement}
\setbeamercolor{block title}{bg=amritaMaroon, fg=white}
\begin{block}{Scientific Premise & Gaps}
\small
\begin{itemize}
\item \textbf{Therapeutic Inequity:} Current clinical frameworks suffer from a severe scarcity of specialists for \textbf{local autism-affected children}, leading to delayed interventions [23].
\item \textbf{Phonetic Complexity:} Standard ASR models exhibit high Word Error Rates (WER) when processing the retroflex consonants and unique prosodic features of \textbf{local autism-affected children} [1, 15, 22].
\item \textbf{Data Vulnerability:} Cloud-based diagnostic pipelines pose significant ethical and privacy risks regarding the biometric voice data of \textbf{local autism-affected children} [25].
\item \textbf{Subjectivity:} Existing progress tracking relies on qualitative clinician observation rather than longitudinal acoustic biomarkers like Vowel Space Area (VSA) [5, 16].
\end{itemize}
\end{block}
\vspace{2mm}
\centering
\footnotesize \textit{Cumulative Gap: The absence of a localized, privacy-preserving, and objective digital scaffolding system for disordered speech [4, 14, 25].}
\end{frame}

% --- NEW SLIDE: OBJECTIVE ---
\begin{frame}{Research Objectives}
\begin{columns}[T]
\begin{column}{0.4\textwidth}
\setbeamercolor{block title}{bg=black, fg=white}
\begin{block}{Technical Objectives}
\scriptsize
\begin{itemize}
\item \textbf{Develop} a quantized, edge-based deep learning pipeline to ensure 100% local data sovereignty [25].
\item \textbf{Integrate} Vocal Tract Length Normalization (VTLN) to adapt to the physiological vocal characteristics of pediatric subjects [19, 21].
\item \textbf{Validate} the system using phonetic benchmarks tailored for regional Dravidian speech patterns [15].
\end{itemize}
\end{block}
\end{column}
\begin{column}{0.5\textwidth}
\setbeamercolor{block title}{bg=amritaMaroon, fg=white}
\begin{block}{Functional Objectives}
\scriptsize
\begin{itemize}
\item \textbf{Evaluate} the efficacy of Reinforcement Learning (RL) in providing adaptive prompt fading for speech motor training [13, 23].
\item \textbf{Establish} a multi-modal biomarker suite (acoustic + visual) to replace subjective clinical scaling [2, 17].
\item \textbf{Bridge} the rural-urban healthcare divide by providing a high-fidelity home therapy tool [9, 24].
\end{itemize}
\end{block}
\end{column}
\end{columns}
\end{frame}

% --- NEW SLIDE: PROOF OF CONCEPT ---
\begin{frame}{Proof of Concept (PoC) & Validation Strategy}
\small
\begin{itemize}
\item \textbf{Sample Inputs:} The system utilizes 16kHz mono-channel disordered audio samples and non-intrusive facial landmark coordinates for stress detection [2, 16].
\item \textbf{Model Behavior:} Employs a hybrid CNN-Transformer architecture optimized via Knowledge Distillation for real-time inference on Raspberry Pi/Jetson Nano [22, 25].
\item \textbf{Adaptivity Logic:} Uses a Deep Deterministic Policy Gradient (DDPG) agent to modulate task difficulty based on the child's real-time success rate [23].
\item \textbf{Output Validation:} System performance is measured against clinician-labeled "Goodness of Pronunciation" (GOP) scores and VSA metrics [5, 8].
\end{itemize}
\vspace{2mm}
\setbeamercolor{block}{bg=gray!10}
\begin{block}{}
\footnotesize \textbf{Feasibility:} Recent studies confirm that 15+ hours of AI-guided interaction leads to measurable improvements in expressive vocabulary [24].
\end{block}
\end{frame}

% --- NEW SLIDE: ARCHITECTURE DIAGRAM ---
\begin{frame}{Architecture Diagram}
    \centering
    \includegraphics[width=\textwidth,height=0.82\textheight,keepaspectratio]{figure-1.png}
\end{frame}

% --- NEW SLIDE: ARCHITECTURE JUSTIFICATION ---
\begin{frame}{Architecture Justification}
\setbeamercolor{block title}{bg=amritaMaroon, fg=white}
\begin{block}{Why Edge-Based Multimodal Design?}
\small
\begin{itemize}
\item \textbf{Privacy by Design:} Processing pediatric data at the edge eliminates the "man-in-the-middle" risks associated with cloud APIs [25].
\item \textbf{Modular Consistency:} A modular pipeline allows independent updates to the Dravidian phonetic model without disrupting the RL scaffolding agent [4, 15].
\item \textbf{Low Latency:} Edge deployment is critical for real-time visual-auditory feedback, which is essential for speech motor learning [12, 18].
\item \textbf{Cost Efficiency:} Utilizing TinyML frameworks reduces the need for expensive high-bandwidth infrastructure in rural settings [22, 25].
\end{itemize}
\end{block}
\end{frame}

% --- NEW SLIDE: COMPARISON WITH EXISTING SYSTEMS ---
\begin{frame}{Comparison with Existing Systems}
    \tiny
    \centering
    \renewcommand{\arraystretch}{1.2}
    \setlength{\tabcolsep}{3pt}
    \begin{tabularx}{\textwidth}{|l|L|L|c|}
        \hline
        \textbf{Parameter} & \textbf{Traditional / Cloud-Based Systems} & \textbf{Proposed Edge-Based System} & \textbf{Citations} \\ \hline
        
        \textbf{Adaptability} & Rigid, predefined lesson paths or manual adjustment by a therapist, leading to potential "prompt dependency." & \textbf{Dynamic Scaffolding:} Utilizes Deep RL (DDPG) to adjust task difficulty in real-time based on phonetic success and child engagement levels. & [13, 18, 23] \\ \hline
        
        \textbf{Multimodality} & Unimodal analysis (Audio only); fails to account for sensory overload or non-verbal cues. & \textbf{Cross-Modal Fusion:} Syncs acoustic biomarkers with non-intrusive visual stress markers and joint attention tracking. & [2, 17] \\ \hline
        
        \textbf{Privacy} & High-risk transmission of sensitive pediatric voice/video data to cloud servers for inference. & \textbf{Data Sovereignty:} 100\% on-device processing via quantized models; uses Federated Learning protocols for privacy. & [25] \\ \hline
        
        \textbf{Autism-Centric} & Designed for general ASR; high rejection rates for atypical prosody, jitter, and shimmer. & \textbf{Disorder-Aware Modeling:} Uses VSA metrics and paralinguistic feature extraction tuned for atypical speech patterns. & [4, 5, 16] \\ \hline
        
        \textbf{Linguistic Focus} & Trained on high-resource Western languages; struggles with Dravidian retroflex phonemes. & \textbf{Regional Normalization:} Implements VTLN and transfer learning specifically for local phonetics in South India. & [15, 19, 22] \\ \hline
        
        \textbf{Scalability} & Requires high-speed internet and expensive recurrent cloud API subscriptions. & \textbf{Offline Resilience:} Operates on low-cost edge hardware, making it feasible for rural/home-based deployment. & [24, 25] \\ \hline
    \end{tabularx}
    
    \vspace{3mm}
    \setbeamercolor{block}{bg=gray!10}
    \begin{block}{}
        \scriptsize \textbf{Research Significance:} While existing tools serve as digital workbooks, our system acts as a \textbf{Real-Time Clinical Scaffolding Agent}. It moves beyond mere transcription by quantifying articulatory motor stability to predict long-term therapeutic outcomes [4, 23, 24].
    \end{block}
\end{frame}
% --- NEW SLIDE: ACADEMIC JUSTIFICATION ---
\begin{frame}{Academic Justification & Collaboration}
\small
\begin{itemize}
\item \textbf{Expert Guidance:} The research direction was significantly refined through valuable inputs from a professor via our project guide, ensuring alignment with current clinical standards [9].
\item \textbf{Collaborative Framework:} This work is carried out in collaboration with domain experts to ensure the RL reward functions mirror actual therapeutic scaffolding techniques [23].
\item \textbf{Methodological Rigor:} Following the precedent of faculty-guided studies, we employ a "Human-in-the-loop" approach for initial model validation [24].
\item \textbf{Peer Precedent:} Collaborative efforts in ASD tech have historically shown higher clinical adoption rates than isolated engineering projects [3, 7].
\end{itemize}
\end{frame}

% --- NEW SLIDE: CONCLUSION & CONFIDENCE ---
\begin{frame}{Conclusion & Confidence Statement}
\setbeamercolor{block title}{bg=amritaMaroon, fg=white}
\begin{block}{Summary of Contributions}
\small
This project provides a technically robust, privacy-first solution for \textbf{local autism-affected children}. By integrating VSA tracking, RL-based adaptivity, and Dravidian-specific phonetic normalization, we bridge a critical gap in assistive healthcare [15, 23, 25].
\end{block}
\vspace{2mm}
\textbf{Faculty Counter-Defense:}
\begin{itemize}
\item \textbf{Feasibility:} Leveraging existing TinyML and Wav2Vec 2.0 architectures ensures the project is achievable within the 8-month timeline [22].
\item \textbf{Readiness:} The team has identified the necessary Dravidian phonetic datasets and edge-hardware constraints for immediate validation [15, 25].
\item \textbf{Significance:} This work aligns with global movements toward decentralized, AI-driven personalized medicine [4, 24].
\end{itemize}
\end{frame}



% --- REFERENCES ---
\section{References}
\begin{frame}[allowframebreaks]{References}
    \tiny
    \begin{enumerate}
        \item J. Zhu, Y. Li, and X. Zhang, "The taste of IPA: Towards open-vocabulary keyword spotting," in Proc. IEEE Int. Conf. Acoustics, Speech, Signal Process. (ICASSP), 2024, pp. 1–5.
        \item A. Landowska, M. Karpus, T. Zawadzka, and M. Robins, "Automatic Emotion Recognition in Children with Autism: A Systematic Review," Sensors, vol. 22, no. 4, p. 1649, 2022.
        \item S. Saeedi, S. Nouri, S. Salari, and H. Kebriaei, "Application of Digital Games for Speech and Language Therapy in Children: A Systematic Review," J. Biomedical Inform., vol. 136, p. 104238, 2022.
        \item M. S. Remya, A. K. Vyas, and K. P. Soman, "Artificial Intelligence Techniques for Speech Disorder Classification: A Review," IEEE Access, vol. 13, pp. 1234–1250, 2025.
        \item S. Sandoval, D. Berisha, R. Utianski, J. Liss, and A. Spanias, "Automatic assessment of vowel space area," J. Acoust. Soc. Am., vol. 134, no. 5, pp. EL477–EL483, 2013.
        \item J. Smith, K. Williams, and P. Brown, "Tell-Your-Day: Augmentative and Alternative Communication for Expressive Sentence Construction," Augment. Altern. Commun., vol. 39, no. 2, pp. 123–135, 2023.
        \item A. Cabibihan, H. Javed, M. Ang Jr., and S. Aljunied, "Robot-Assisted Joint Attention and Verbal Initiation Training for Children with ASD," IEEE Trans. Neural Syst. Rehabil. Eng., vol. 30, pp. 1845–1855, 2022.
        \item L. Chen, R. Zhang, and Y. Liu, "AI-Based Real-Time Speech Feedback Using ASR and Automated Scoring," IEEE Trans. Learning Technol., vol. 17, no. 3, pp. 234–245, 2024.
        \item M. Johnson, D. Peterson, and K. Roberts, "Traditional and Innovative Speech Therapy Approaches for Children with Autism Spectrum Disorder," J. Autism Dev. Disord., vol. 53, no. 6, pp. 2456–2470, 2023.
        \item R. Patel, S. Connaghan, D. Franco, E. Edsall, D. Forgit, and L. Olsen, "The Caterpillar: Automatic Speech Recognition for Children with Autism Spectrum Disorder," J. Speech Lang. Hear. Res., vol. 65, no. 11, pp. 4321–4335, 2022.
        \item T. Garcia, L. Rodriguez, and M. Fernandez, "Immersive Virtual Reality Training for Practicing Social Scripts in Adolescents with ASD," Virtual Reality, vol. 24, no. 3, pp. 401–412, 2020.
        \item H. Zhao, W. Chen, and X. Wang, "Real-time Visual Feedback Systems for Pitch and Volume Control in Speech Therapy," IEEE Trans. Biomed. Eng., vol. 68, no. 9, pp. 2701–2710, 2021.
        \item Y. Chen, T. Wu, and H. Lin, "Self-Learning Bot using Transformers for Intent Decoding in Pediatric Speech Therapy," in Proc. Conf. Empirical Methods Natural Lang. Process. (EMNLP), 2022, pp. 5678–5689.
        \item M. White, P. Green, and S. Taylor, "Mobile Applications for Articulation Training and Adaptive Phonetic Profiling," Int. J. Speech-Lang. Pathol., vol. 21, no. 4, pp. 378–389, 2019.
        \item S. Lee, J. Park, and K. Kim, "Knowledge Transfer and Meta-learning for Dravidian Language Adaptation in Speech Recognition," IEEE/ACM Trans. Audio, Speech, Language Process., vol. 31, pp. 1789–1802, 2023.
        \item P. Nguyen, D. Le, and M. Tran, "Atypical Prosody in Autism: Quantification Using openSMILE Toolkit," J. Autism Dev. Disord., vol. 51, no. 8, pp. 2801–2815, 2021.
        \item F. Wang, X. Li, and Y. Zhao, "Multi-modal Fusion: Synchronizing Vocalization and Gesture Tracking for Communication Assessment," IEEE Trans. Affect. Comput., vol. 13, no. 4, pp. 1891–1903, 2022.
        \item D. Miller, A. Johnson, and B. Davis, "Feedback Timing in Motor Learning: Impact of Immediate versus Delayed Responses," Motor Control, vol. 24, no. 2, pp. 201–215, 2020.
        \item G. Kim, S. Choi, and J. Lee, "Acoustic Modeling for Disordered Speech: VTLN and Generative Adversarial Networks," in Proc. Interspeech, 2018, pp. 3267–3271.
        \item K. Sato, Y. Tanaka, and H. Nakamura, "Electropalatography Biofeedback: Visualizing Tongue-Palate Contact for Articulation Training," Clin. Linguist. Phonet., vol. 31, no. 5, pp. 345–360, 2017.
        \item P. Govalkar, J. Fischer, F. Zalkow, and C. Dittmar, "A Comparison of Recent Neural Vocoders for Speech Signal Reconstruction," in Proc. Int. Conf. Digital Audio Effects (DAFx), 2019, pp. 1–8.
        \item N. D. Al Futaisi, Z. Alkatheeri, and A. Al Hashmi, "The Noor Project: Wav2Vec 2.0 Transformer for Arabic ASD Speech Classification," IEEE Access, vol. 13, pp. 12456–12468, 2025.
        \item Y. Chen, H. Wang, and L. Zhang, "Deep Deterministic Policy Gradient for Autism Screening and Adaptive Intervention Strategies," IEEE J. Biomed. Health Inform., vol. 29, no. 1, pp. 234–246, 2025.
        \item C.-H. Chueh, M.-L. Wu, and Y.-T. Chen, "Generative AI-Powered Clinical Language Therapy: A Randomized Controlled Trial with 2000+ Home Recordings," JAMA Pediatrics, vol. 179, no. 3, pp. 287–295, 2025.
        \item N. Pavlidis, I. Mavridis, and E. Georgiadis, "Federated Learning for Privacy-Preserving Anomaly Detection in Autism Spectrum Disorder Diagnosis," IEEE Trans. Inform. Forensics Security, vol. 19, pp. 4567–4580, 2024.
    \end{enumerate}
\end{frame}

\begin{frame}
    \centering    
    {\Huge \textbf{\textcolor{amritaMaroon}{Thank You}}}
\end{frame}

\end{document}
